\section{Using the package}
\label{sec:usage}

The package \citep{cvxcla} is installed with \code{pip install cvxcla}; it is pure
\proglang{Python} and needs \pkg{NumPy}, \pkg{SciPy} and \pkg{cvx-linalg}
(Computational details). The workflow goes through a single class, \code{CLA}. The user supplies $\mu$, the
covariance (a dense array or any backend of \S\ref{sec:operators}), the bounds, and
the constraints; constructing the object traces the whole frontier. The example traces
a small long-only, fully-invested problem and reads off the maximum-Sharpe portfolio.

\begin{CodeInput}
import numpy as np
from cvxcla import CLA, FactorCovariance

rng = np.random.default_rng(1)
n, k = 6, 2
d = rng.uniform(0.1, 0.5, n)            # idiosyncratic variances
U = rng.standard_normal((n, k))         # factor loadings
Delta = np.diag(rng.uniform(0.5, 2.0, k))
Sigma = np.diag(d) + U @ Delta @ U.T    # a dense 2-factor covariance
mean = rng.standard_normal(n)

cla = CLA(
    mean=mean, covariance=Sigma,
    lower_bounds=np.zeros(n), upper_bounds=np.ones(n),
    a=np.ones((1, n)), b=np.ones(1),    # fully invested: sum(w) = 1
)
print("turning points:", len(cla.turning_points))
sr, w = cla.frontier.max_sharpe
print("max Sharpe:", round(sr, 4))
print("weights:", np.round(w, 3))
\end{CodeInput}

\begin{CodeOutput}
turning points: 7
max Sharpe: 2.4037
weights: [0.    0.    0.704 0.296 0.    0.   ]
\end{CodeOutput}

Switching to a structured backend changes only the covariance argument; because the
factor backend represents the \emph{same} $\Sig$, it returns the same frontier to round-off.

\begin{CodeInput}
factor = FactorCovariance(d=d, u=U, delta=np.diag(Delta))
cla_factor = CLA(
    mean=mean, covariance=factor,
    lower_bounds=np.zeros(n), upper_bounds=np.ones(n),
    a=np.ones((1, n)), b=np.ones(1),
)
w_dense = np.array([t.weights for t in cla.turning_points])
w_factor = np.array([t.weights for t in cla_factor.turning_points])
print("identical frontier:", np.allclose(w_dense, w_factor))
\end{CodeInput}

\begin{CodeOutput}
identical frontier: True
\end{CodeOutput}

A fluent builder, \code{CLA.problem}, assembles the same arguments step by step through
methods such as \code{.budget}, \code{.equality} and \code{.inequality}; an extra \code{.equality} row,
such as a fixed sector weight, makes the first turning point the linear program of
\S\ref{sec:first}.


\subsection{The interface}
\label{sec:api}

\code{CLA} takes the problem as keyword arguments, all \pkg{NumPy} arrays of floats, and
traces the frontier in its constructor. Required are \code{mean} ($\mu$, length $n$),
\code{covariance} (an $n\times n$ array, which is wrapped as a \code{DenseCovariance},
or any backend of \S\ref{sec:operators}), \code{lower\_bounds} and \code{upper\_bounds}
(length $n$), and \code{a}, \code{b} ($m\times n$ and length $m$; a budget is
\code{a = np.ones((1, n))}, \code{b = np.ones(1)}). Optional are \code{g}, \code{h}
($p\times n$ and length $p$) for $Gw \le h$, \code{leverage} for the cap
$\lVert w\rVert_1 \le c$, and \code{tol} (default $10^{-5}$), the tolerance for deciding
that a weight sits on a bound and for validating each turning point; the event ordering
uses its own relative window of at most $10^{-10}$ (\S\ref{sec:events}). The arrays (and a
dense covariance) are copied into read-only arrays when the object is built, so
changing the caller's arrays afterwards cannot put the stored turning points out of step
with the data; a backend is used as given. Non-finite data are refused at this point,
as are NaN bounds. Redundant equality rows are
dropped, and \code{cla.a}, \code{cla.b} then hold the reduced system.

The result is \code{cla.turning\_points}, a list of \code{TurningPoint} objects with the
\code{weights}, the value \code{lamb} of $\lambda$, the boolean mask \code{free} of free
assets, and the mask \code{active\_ineq} of active rows of $G$. The maximum-return
vertex appears twice, at $\lambda = \infty$ and again at the first finite turning point,
where the path leaves it; \code{len(cla)} is the number of turning points.
\code{cla.frontier} wraps them as a \code{Frontier}
(\srcfile{src/cvxcla/types.py}{types.py}), which exposes the \code{weights},
\code{returns}, \code{variance}, \code{volatility} and \code{sharpe\_ratio} of the
turning points and a \code{plot} method. Because the frontier is piecewise linear in
weight space, \code{interpolate(num)} returns intermediate portfolios with no
discretisation error.

\subsection{The frontier and the maximum-Sharpe portfolio}
\label{sec:frontier}

\code{max\_sharpe} maximises $\mu\trans w/\sqrt{w\trans\Sig w}$, the Sharpe ratio with a
zero risk-free rate (supply excess returns as $\mu$ otherwise), over the computed
frontier. On a segment $w(t) = (1-t)\,w_0 + t\,w_1$, $t \in [0,1]$, the return is affine
and the variance quadratic in $t$,
\begin{equation}
S(t) = \frac{a_0 + a_1 t}{\sqrt{c_0 + c_1 t + c_2 t^2}}, \qquad
\begin{aligned}
a_0 &= \mu\trans w_0, & c_0 &= w_0\trans\Sig w_0, & c_2 &= \Delta\trans\Sig\Delta,\\
a_1 &= \mu\trans\Delta, & c_1 &= 2\,w_0\trans\Sig\Delta, & \Delta &= w_1 - w_0,
\end{aligned}
\end{equation}
and the numerator of $S'(t)$ is affine in $t$ (its $t^2$ terms cancel), so $S$ has at
most one stationary point,
\begin{equation}\label{eq:tstar}
t^\star = \frac{a_0 c_1 - 2 a_1 c_0}{a_1 c_1 - 2 a_0 c_2}.
\end{equation}
The maximum on the segment is the largest of $S(0)$, $S(1)$ and, when
$0 < t^\star < 1$, $S(t^\star)$; a stationary point outside the segment is ignored, an
optimum at a turning point is found as an end point, and when $a_1 c_1 = 2 a_0 c_2$ there
is no stationary point and only the end points are compared. The code applies this
rule to every segment of the frontier and returns the best candidate, so the result is
the exact maximum over the computed frontier, without relying on the ratio being
unimodal, at the cost of one matrix--vector product with $\Sig$ per segment.

Points of zero variance are skipped (\code{sharpe\_ratio} reports NaN there); with no
positive variance \code{max\_sharpe} raises a \code{ValueError}, and with no positive
expected return it warns and returns the least negative ratio. Whether the ratio,
computed from estimated $\mu$ and $\Sig$, predicts future performance is a statistical
question outside the package (\S\ref{sec:estimators}).

\subsection{Supported inputs and errors}
\label{sec:errors}




Table~\ref{tab:errors} sorts the inputs into kinds, each with its own exception and
remedy, and each row is a test in the package
(\srcfile{tests/test_errors.py}{tests/test\_errors.py} and
\srcfile{tests/test_degenerate_vertex.py}{tests/test\_degenerate\_vertex.py}). All are raised by the constructor. The package errors derive from \code{CLAError} and also from the built-in
exception raised before (\code{ValueError}, or \code{RuntimeError}), so existing
handlers keep working.



A sample covariance from fewer observations than assets is a case that used to be
declined and no longer is. It is singular, and its free block becomes singular partway
along the frontier. The budget row still holds the free assets on its null space, where
the block stays positive definite, so the trace goes on there (Appendix~\ref{sec:degeneracy}).
A ridged Gram backend (\S\ref{sec:gram}) traces a different, regularised frontier:

\noindent\begin{minipage}{\textwidth}
\begin{CodeInput}
from cvxcla import GramCovariance

returns = 0.01 * rng.standard_normal((4, n)) + 0.001    # 4 days, 6 assets: rank 3
long_only = dict(lower_bounds=np.zeros(n), upper_bounds=np.ones(n),
                 a=np.ones((1, n)), b=np.ones(1))
cla = CLA(mean=returns.mean(axis=0),
          covariance=np.cov(returns, rowvar=False), **long_only)
ridged = CLA(mean=returns.mean(axis=0),
             covariance=GramCovariance(returns, ridge=1e-4), **long_only)
print("turning points:", len(cla), "| with a ridge:", len(ridged))
\end{CodeInput}
\begin{CodeOutput}
turning points: 5 | with a ridge: 7
\end{CodeOutput}
\end{minipage}

Release~2.3.3 declined the first of these. Every turning point of the trace above
matches an interior-point solve of the same quadratic program.

\begin{table}[tb]
\centering
\small
\renewcommand{\arraystretch}{1.25}
\newcolumntype{L}[1]{>{\raggedright\arraybackslash}p{#1}}
\begin{tabular}{@{}L{0.62\textwidth}L{0.32\textwidth}@{}}
\toprule
Input & Remedy \\
\midrule
\multicolumn{2}{@{}l}{\textbf{Malformed input}: \code{ValueError}, before the walk} \\
Non-symmetric $\Sig$; non-finite data or NaN bounds; an indefinite factor $\Delta$;
wrong shapes & correct the data \\
\midrule
\multicolumn{2}{@{}l}{\textbf{Infeasible}: \code{InfeasibleProblemError}, before the walk} \\
Crossed bounds; an unreachable budget; constraints with no feasible point; dependent
equality rows that contradict the others & change the constraints \\
\midrule
\multicolumn{2}{@{}l}{\textbf{Harmless degeneracy}: traced, nothing raised} \\
Redundant equality rows; duplicated, never-binding or locally redundant inequality
rows; a degenerate maximum-return vertex whose optimal dual set has a vertex, and tied
maximum returns (\S\ref{sec:first}); a singular covariance whose free blocks the active
rows keep positive definite (Appendix~\ref{sec:degeneracy}) & -- \\
\midrule
\multicolumn{2}{@{}l}{\textbf{Unsupported degeneracy}: \code{DegenerateProblemError}} \\
Before the walk: a degenerate vertex whose optimal dual set has no vertex, as when a
tight row touches only assets fixed at both bounds; infinite bounds that make the
return unbounded & perturb the bounds or constraints \\
During the walk: a free block singular even on the null space of the active rows, or
active rows dependent on the free set & a ridge or factor model; loosen a cap \\
\midrule
\multicolumn{2}{@{}l}{\textbf{Numerical failure}: \code{NumericalError}, during the walk} \\
Iteration cap; projection fails; a turning point violates a constraint
& regularise $\Sig$ or relax \code{tol}; report it \\
\bottomrule
\end{tabular}
\caption{Supported and unsupported inputs, grouped by the exception each raises and
when, with the usual remedy. ``Before the walk'' is input validation and the first
vertex; ``during the walk'' is a turning point of the trace.}
\label{tab:errors}
\end{table}

